\documentclass[../../../uem_mat_tok_q.tex]{subfiles}

\begin{document}
    平面上で $\mathrm{AB}=\mathrm{AC}=1$ である二等辺三角形ABCを考える．%
    正の実数 $r$ に対し，A，B，Cそれぞれを中心とする半径 $r$ の円3つを合わせた領域を %
    $D_r$ とする．ただし，この問いでは，三角形と円は周とその内部からなるものとする．%
    辺AB，AC，BCがすべて $D_r$ に含まれるような最小の $r$ を $s$，%
    三角形ABCが $D_r$ に含まれるような最小の $r$ を $t$ と表す．

    \begin{enumerate}
        \item $\angle\mathrm{BAC}=\myfraca{\pi}{3}$ のとき，$s$ と $t$ を求めよ．
        \item $\angle\mathrm{BAC}=\myfraca{2\pi}{3}$ のとき，$s$ と $t$ を求めよ．
        \item $0<\theta<\pi$ を満たす $\theta$ に対して，%
        $\angle\mathrm{BAC}=\theta$ のとき，$s$ と $t$ を $\theta$ を用いて表せ．
    \end{enumerate}
\end{document}